\section{MMU and Virtual Memory}

\subsection{The problem: software should not always use physical addresses}

Without an MMU, every address produced by software is a physical address. If code executes:

\begin{lstlisting}[language={}]
sw x5, 0(x10)
\end{lstlisting}

and \texttt{x10} contains \texttt{0x10000000}, then the CPU writes physical address \texttt{0x10000000} directly. That address might be RAM, a UART register, a timer register, another program’s memory, or an invalid area.

This is simple, but it has problems:

\begin{itemize}
\item programs can corrupt each other;
\item user software can access device registers directly;
\item each program must know where it is physically loaded;
\item operating systems cannot easily isolate processes;
\item memory can be fragmented and still needs to appear contiguous to software.
\end{itemize}
\subsection{Virtual address vs physical address}

Virtual memory introduces a translation step:

\begin{lstlisting}[language={}]
software address  ->  MMU translation  ->  physical memory/bus address
virtual address       page table walk      physical address
\end{lstlisting}

The program produces a \textbf{virtual address}. The memory system needs a \textbf{physical address}. The MMU translates between them.

This lets the operating system give each program its own address space. Two programs can both think they use address \texttt{0x00001000}, but the MMU can translate those virtual addresses to different physical pages.

\subsection{What a page is}

A page is a fixed-size block of virtual memory. In RISC-V Sv32, the normal page size is 4 KiB. This means the lower 12 bits of an address are the offset within the page:

\begin{lstlisting}[language={}]
4 KiB = 4096 bytes = 2^12 bytes
\end{lstlisting}

So address bits \texttt{[11:0]} are not translated. They are copied into the physical address. The upper bits select which page is being accessed.

\subsection{What a page table is}

A page table is a data structure in memory that tells the MMU how to translate virtual pages into physical pages. Each entry is called a PTE: Page Table Entry.

A PTE contains:

\begin{itemize}
\item a physical page number;
\item valid/read/write/execute/user permission bits;
\item other status bits.
\end{itemize}
The MMU reads PTEs from memory during a page-table walk.

\subsection{Why multi-level page tables exist}

A flat page table for a 32-bit address space could be large. Multi-level page tables solve this by splitting the virtual page number into pieces. The first piece indexes a root table. That entry either points to a second-level table or directly maps a large page. The second piece indexes the second-level table.

This saves memory because page tables only need to exist for regions that are actually mapped.

\subsection{Now Sv32: the RISC-V 32-bit virtual memory scheme}

Only after understanding pages and page tables does Sv32 make sense.

Sv32 is the RISC-V virtual memory scheme for 32-bit virtual addresses. It divides the address like this:

\begin{lstlisting}[language={}]
31          22 21          12 11           0
+-------------+--------------+--------------+
|   VPN[1]    |    VPN[0]    | page offset  |
|   10 bits   |    10 bits   |   12 bits    |
+-------------+--------------+--------------+
\end{lstlisting}

\begin{itemize}
\item \texttt{VPN[1]} indexes the level-1/root page table.
\item \texttt{VPN[0]} indexes the level-0 page table.
\item \texttt{page offset} is copied to the physical address.
\end{itemize}
\subsection{13.7 \texttt{satp}: where translation begins}

The CPU needs to know where the root page table is. RISC-V stores that information in the \texttt{satp} CSR.

For RV32 Sv32, the useful layout is:

\begin{lstlisting}[language={}]
31      30          22 21                         0
+--------+-------------+---------------------------+
| MODE   |    ASID     |          root PPN         |
| 1 bit  |   9 bits    |          22 bits          |
+--------+-------------+---------------------------+
\end{lstlisting}

In this implementation:

\begin{itemize}
\item \texttt{satp[31] = 1} enables Sv32 translation;
\item \texttt{satp[21:0]} is used as the root page table physical page number;
\item ASID is structurally present in the format but not used for a TLB because this educational MMU has no TLB.
\end{itemize}
\subsection{PTE format and permissions}

\texttt{rtl/mmu.v} uses these PTE bits:

\begin{longtable}{|>{\RaggedRight\arraybackslash}p{0.28\textwidth}|>{\RaggedRight\arraybackslash}p{0.32\textwidth}|>{\RaggedRight\arraybackslash}p{0.30\textwidth}|}
\hline
\textbf{Bit} & \textbf{Name} & \textbf{Meaning} \\
\hline
\endfirsthead
\hline
\textbf{Bit} & \textbf{Name} & \textbf{Meaning} \\
\hline
\endhead
0 & V & Valid \\
\hline
1 & R & Readable \\
\hline
2 & W & Writable \\
\hline
3 & X & Executable \\
\hline
4 & U & User-accessible \\
\hline
5 & G & Global, noted but not central here \\
\hline
6 & A & Accessed, noted but not updated here \\
\hline
7 & D & Dirty, noted but not updated here \\
\hline
31:10 & PPN & Physical page number \\
\hline
\end{longtable}

The implementation treats a PTE as invalid if \texttt{V=0} or if \texttt{R=0} and \texttt{W=1}, which is an invalid combination.

\subsection{Page-table walk}

For a data access in user mode when Sv32 is enabled, the MMU performs:

\begin{lstlisting}[language={}]
pte1_addr = satp.root_ppn * 4096 + VPN[1] * 4
pte1 = memory[pte1_addr]

if pte1 is a leaf:
    map as a 4 MiB superpage
else:
    pte0_addr = pte1.ppn * 4096 + VPN[0] * 4
    pte0 = memory[pte0_addr]
    pte0 must be a leaf mapping a 4 KiB page
\end{lstlisting}

A leaf is detected when \texttt{R} or \texttt{X} is set. A level-1 leaf maps a superpage. A level-0 leaf maps a normal 4 KiB page.

\subsection{How physical addresses are formed}

For a 4 KiB page:

\begin{lstlisting}[language={}]
physical_address = pte0.PPN || virtual_address[11:0]
\end{lstlisting}

For a 4 MiB superpage:

\begin{lstlisting}[language={}]
physical_address = pte1.PPN[1] || virtual_address[21:0]
\end{lstlisting}

\texttt{rtl/mmu.v} implements these as \texttt{pa\_4k} and \texttt{pa\_super}.

\subsection{When translation is active}

In \texttt{rtl/mmu.v}, translation is active only when:

\begin{lstlisting}[language={}]
priv == user mode AND satp[31] == 1
\end{lstlisting}

Machine mode bypasses translation and uses addresses directly. This keeps machine-mode trap handlers and boot code simple.

\subsection{Page faults}

The MMU reports a page fault when:

\begin{itemize}
\item a PTE is invalid;
\item a required second-level leaf is missing;
\item user mode accesses a page without \texttt{U=1};
\item a load is attempted without read permission;
\item a store is attempted without write permission.
\end{itemize}
The CPU converts this into RISC-V trap causes:

\begin{longtable}{|>{\RaggedRight\arraybackslash}p{0.38\textwidth}|>{\RaggedRight\arraybackslash}p{0.52\textwidth}|}
\hline
\textbf{Cause} & \textbf{Meaning} \\
\hline
\endfirsthead
\hline
\textbf{Cause} & \textbf{Meaning} \\
\hline
\endhead
13 & Load page fault \\
\hline
15 & Store/AMO page fault \\
\hline
\end{longtable}

\subsection{Two MMU implementations in this project}

There are two MMU styles:

\subsubsection{\texttt{rtl/mmu.v} — combinational educational MMU}

This module has dedicated page-table read ports:

\begin{lstlisting}[language={}]
walk_addr1 / walk_data1
walk_addr2 / walk_data2
\end{lstlisting}

That lets the page walk complete combinationally in one cycle. This is easy to understand and useful for simulation, but unrealistic for a single-port block RAM system.

\subsubsection{\texttt{rtl/cpu\_mc\_mmu.v} — synthesizable multi-cycle walker}

This version sequences the walk through states such as:

\begin{lstlisting}[language={}]
PTW1  -> drive level-1 PTE address
PTW1D -> receive/check level-1 PTE
PTW0D -> receive/check level-0 PTE
ACC   -> perform translated memory access
ACCD  -> receive load data
PF    -> raise page-fault trap
\end{lstlisting}

This is closer to real hardware because the page-table walker shares the memory/data port and waits for registered memory data.

\subsection{What is simplified}

This MMU is educational. It does not implement a TLB, does not update Accessed/Dirty bits, and focuses on data-side translation. A production MMU would add instruction-side translation, TLBs, permission details such as SUM/MXR, page-fault priority details, and more complete privilege behavior.

\bigskip
